\section{背景} \label{sec:background}
聴音トレーニング（Ear Training）とは，聴覚を用いて音高，音程，旋律，和音，リズムといった音楽的要素を識別する技術を習得するための訓練である．これは音楽家の育成において極めて重要な役割を果たしており，その最終的な到達点は，聴覚のみを用いて音楽を正確に表現し，伝達する「インナーヒアリング」の獲得にある\cite{woody2012playing}．この文献で，多様な形態が提案されているが，こうした聴覚能力の獲得における基礎は旋律の認識であり，そのためには構成音間の「音程」を区別し認識する能力を重要視している．したがって，音程認識能力の向上は，音楽学習者が優れたインナーヒアリングを獲得するための最も基礎的な手段であるといえる．

従来，聴音トレーニングは指導者やパートナーがピアノ等で課題を演奏し，学習者の回答を評価するという人的リソースに依存した形式で行われてきた．しかし，計算機技術の発展に伴い，様々な訓練用ソフトウェアが開発されている．その歴史は古く，1970年代半ばにはPLATOメインフレーム上で動作する「G.U.I.D.O.」\footnote{https://www.learntechlib.org/p/168759/}が開発され，大学の音楽学生向けに音程・旋律・和音等の認識指導が行われた．また，1996年にリリースされた「EarMaster」\footnote{https://www.earmaster.jp/}は，現在でも広く利用されている代表的なソフトウェアであり，音程比較からリズム認識まで多岐にわたるトレーニング機能を提供している．これらはスタンドアロン，Webブラウザ，モバイルデバイスなど多様なプラットフォームで利用可能であり，独習の機会を広げている．

さらに，近年急速に普及しているVirtual Reality（VR）技術は，視覚，聴覚，触覚を含む多感覚入力を統合した仮想環境（Virtual Environments）を提供する技術である．SerafinらによりこうしたVRおよびAR（拡張現実）の技術とソフトウェアについて音楽教育分野に関する包括的なレビューが行われている\cite{serafin2017considerations}．また現在VRでの音楽教育分野における研究として特定の楽器における楽器支援\cite{gao2024application}や、既存の音楽レッスンによる拡張\cite{liu2025research}が挙げられ，本研究は後者に位置づけられる．こうした研究がVRを扱う利点として高価な機器の準備や教師不足が挙げられている．またVRは既存のウェブアプリケーションよりも空間的要素を高めることができ，音楽学習者が楽器の演奏や視聴における近しい体験が提供できる利点もある．また現在の聴音トレーニングにおけるインターフェースの多くはピアノの鍵盤のような「水平・平面的な配置」に依存している．PedersonらはVRでの聴音トレーニングにおいて、ステレオ音響を音程認識の手がかりとして影響を与えていることが示唆されている\cite{pedersen2020spatialized}。これはステレオ音響がVR技術での学習において有益である可能性について期待できる．従来の平面的なインターフェースに対し，音のピッチ（高低）は空間的な高低に対して結びついていることが示されており\cite{rusconi2006spatial}，VR空間であれば，物理的な制約を受けずに，この「音の高さ」と「空間的な高さ」を一致させた立体的な音源配置が可能である．

そこで本研究では，VR技術を用いた聴音トレーニングについて，既存手法\cite{pedersen2020spatialized, fletcher2019virtual}である平面的インターフェースを模した前方半円状配置と，音程の周波数を空間の距離に対応させた円状配置したシステムにおける聴音トレーニングを比較検討することで，音楽教育分野のさらなる発展に向けた可能性を模索する．

\section{研究目的}
本研究の目的は，VR空間内の聴音トレーニングにおいて，音源を立体的に配置することが学習者の音程認識に与える影響を明らかにすることである．平面的なインターフェースとして前方半円状配置と，音程の周波数に基づく音源の高さの調整を行った円状配置の2つのシステムを用いた．18名を実験参加者とし，それぞれの音楽経験を考慮した上での，2つの音が順になり「ピッチインターバル（音程識別」を聞き分けるテストを行い，さらなる評価を行っている．本研究の成果は，将来的なVR音楽教育システムの設計において，更なる発展に寄与することを目指す．

\section{論文構成}
本論文の構成は以下の通りである．\\

第1章　序論

　本研究の背景および研究目的について述べる．

第2章　関連研究

　本研究に関連する研究について述べる．\\

第3章　設計

　実験に使用したゲームの設計について述べる．

第4章　実験

　実施した評価実験について述べる．

第5章　結果

　評価実験にて得られた結果について示す．

第6章　考察

　結果に基づいた考察を述べる．

第7章　結論

　本研究の結論と将来課題について述べる．